\section{The \texttt{Source\_lab} McXtrace Component}
Laboratory x-ray source.

\subsection*{Identification}
\begin{itemize}
  \item \textbf{Author:} Erik B. Knudsen
  \item \textbf{Origin:} Kgs. Lyngby
  \item \textbf{Date:} May 2012
\end{itemize}

\subsection*{Description}
Model of a laboratory x-ray tube, generating x-rays by bombarding a target by electrons. Given a input energy E0 of the electron beam, x-rays are emitted from the accessible emission lines The geometry of the tube is assumed to be: \# The electron beam hits a slab of surface material surface at a right angle illuminating an area of width by height, \# where width is measured along the component X-axis. \# The centre of the electron beam at the anode surface is the origin of the component. \# The Z-axis of the component points at the centre of the exit window (focus\_xw by focus yh) placed at a distance dist from the origin. \# The angle between the Z-axis and the anode surface is the take\_off angle. For a detailed sketch of the geometry see the componnent manual.

The Bremsstrahlung emitted is modelled using the model of Kramer (1923) as restated in International Tables of Crystallography C 4.1 Characteristic radiation is modelled by Lorentzian (default) or Gaussian energy profiles with line-energies from Bearden (1967), widths from Krause (1979) and intensity from Honkimäki (1990) and x-ray data booklet. Absoprtion of emitted x-rays while travelling through the target anode is included.

Example: Source\_lab(material\_datafile="Cu.txt",Emin=1, E0=80)

\subsection*{Input parameters}
Parameters in \textbf{boldface} are required; the others are optional.

\begin{longtable}{p{0.22\textwidth}p{0.12\textwidth}p{0.46\textwidth}p{0.14\textwidth}}
\toprule
\textbf{Name} & \textbf{Unit} & \textbf{Description} & \textbf{Default} \\
\midrule
\endhead
material\_datafile & string & Name of datafile which describes the target material. & "Cu.txt" \\
width & m & Width of electron beam impinging on the anode. & 1e-3 \\
height & m & Height of electron beam impinging on the anode. & 1e-3 \\
thickness & m & Thickness of the anode material slab. & 100e-6 \\
E0 & kV & Acceleration voltage of xray tube. & 20 \\
Emax & keV & Maximum energy to sample. Default (Emax=0) is to set it to E0. & 0 \\
Emin & keV & Minimum energy to sample. & 1 \\
focus\_xw & m & Width of exit window. & 5e-3 \\
focus\_yh & m & Height of exit window. & 5e-3 \\
take\_off & deg & Take off angle of beam centre. & 6 \\
dist & m & Distance between centre of illuminated target and exit window. & 1 \\
tube\_current & A & Electron beam current. & 1e-3 \\
frac & 0-1 & Fraction of statistic to use for Bremsstrahlung. & 0.1 \\
lorentzian & 0/1 & If nonzero Lorentzian (more correct) line profiles are used. & 1 \\
xwidth & m & Width of the anode material slab. & 0 \\
yheight & m & Height of the anode material slab. & 0 \\
exit\_window\_refpt & m & If set, the AT position and exit window will coincide (legacy behaviour). & 0 \\
\bottomrule
\end{longtable}

\subsection*{Links}
\begin{itemize}
  \item Component source code found in file \texttt{Source\_lab.comp}.
\end{itemize}
\IfFileExists{sources/Source_lab_static.tex}{\input{sources/Source_lab_static.tex}}{}